\section{Conclusion\label{sec:conclusion}}

This study examines how code text, static dependencies, folder organization, and member-level information contribute to code-to-architecture mapping on ten systems. The findings show that mapping performance depends not only on which information is available but also on how it is represented and used.

Representations that encode dependency information generally outperform the evaluated lexical and zero-shot code representations. Node2Vec and dependency-guided fine-tuned UniXcoder outperform Word2Vec on most systems, and fine-tuning improves UniXcoder over its zero-shot counterpart on nearly all systems. However, neither Node2Vec nor fine-tuned UniXcoder is consistently best across systems.

Folder nodes produce the largest and most consistent gains within the HGT variants, but a directory rule using the same labeled files reaches similar accuracy on most systems, and folder paths supplied as file features match or exceed folder nodes on most systems. These gains do not consistently translate into better predictions on files the rule assigns incorrectly. Member nodes provide no consistent improvement with one or two HGT layers, and their effects cannot generally be separated from changes in feature construction and placement. These structural findings are specific to HGT; whether they persist with other predictors remains untested.

Keeping the input representations fixed while enabling and disabling access to dependency-neighborhood information shows that neighbor aggregation can complement representations that already encode dependencies. With undirected aggregation, the median change across systems is positive on both metrics for every predictor and representation. In this setting, GAT achieves the highest median macro-F1, while HGT gains the least. The better aggregation direction varies between systems, although the experiments do not isolate propagation direction from neighborhood size. Evaluation should therefore distinguish information already encoded in the representations from the contribution of neighborhood access. The findings support assessing these contributions for the intended mapping task, rather than treating any representation or predictor as generally preferable.

\parg{Future work.} Future research could clarify how different aspects of graph context contribute to the observed differences. Comparing incoming, outgoing, and combined neighborhoods at matched sizes would help distinguish the effects of propagation direction and neighborhood size. Adding member nodes while preserving each file's own features would help separate the effects of graph structure and feature placement. Another direction is configuration selection on unfamiliar systems. Homophily, which helped interpret the system differences, requires reference module assignments, so selection rules based on properties observable in the dependency graph alone could be developed on some systems and evaluated on others. Finally, a model trained on one version could be evaluated on new and changed files in later versions without retraining, to test whether the learned information supports mapping as systems evolve.
